% RESEARCH_DRAFT_SHA256: f4d66b0ecc25fbac63aa65d8d519c7590ad409af3b06b3e63dad03d799eb23aa
% Layout and source-reference wrappers expanded for the web renderer.
\documentclass[11pt]{article}
\usepackage[margin=25mm]{geometry}
\usepackage{fontspec,unicode-math}
\setmainfont{Latin Modern Roman}
\setmathfont{Latin Modern Math}
\usepackage{amsmath,amssymb,mathtools,xurl,hyperref}
\hypersetup{colorlinks=true,urlcolor=blue,linkcolor=blue}
\setcounter{secnumdepth}{0}
\setlength{\parindent}{0pt}
\setlength{\parskip}{5pt}
\setlength{\emergencystretch}{3em}
\newcommand{\sref}[2]{\hyperlink{source:#1:#2}{[S#2]}}
\newcommand{\cataloguescope}[1]{\paragraph{Source record.}#1}
\newcommand{\floor}[1]{\left\lfloor#1\right\rfloor}
\begin{document}
\section*{Research-attempt provenance}
This standalone research attempt retains the complete original definition body below,
including its historical source and preparation notes. Those notes describe the
earlier catalogue preparation; the new attempt and its verification are identified
separately after that body. No original ID, catalogue status, or source statement is changed.

\section{A certifiable bridge for intermediate PreTeXt users}\label{statement}

\subsection{Definitions and mathematical statement}

\paragraph{Shared notation.}Standard number systems are $\mathbb N,\mathbb Z,\mathbb Q,\mathbb R,\mathbb C$. Specialized notation and index ranges follow the cited source and the definitions below.


Reference documentation describes a system's vocabulary, options, and behavior for users seeking specific information. An introductory tutorial teaches a path through the system. An intermediate user has some prior familiarity, but the prerequisite knowledge must be stated to organize suitable documentation.
\paragraph{Record scope and qualifications.}
The catalogue title is retained as metadata. Named certificates, contracts, bounds or frameworks in that title are not asserted as results of the original source extract.
The reconstructed extract does not identify the particular guide section or an explicit set of user competencies.

Catalogue screening limitation: PreTeXt software design, documentation, or publication feature request; no mathematical problem is specified by the supplied passage.

The source snapshot labels this record partially solved; this is catalogue metadata rather than independent verification of a mathematical result or the current completion of a project.
The source text below is rendered literally where necessary; damaged formulas and composite extracts are not silently repaired into a new mathematical claim.
The supplied catalogue marks this formulation as an unverified reconstruction; missing assumptions are not certified by the source link.
\paragraph{Statement recorded in the supplied source.}
Reference for intermediate users: what can go here? \href{https://github.com/PreTeXtBook/community/wiki/Moderated-Problem-Session}{[S2]}
\paragraph{Source status.}
The supplied catalogue labels this source record partially solved; that metadata was checked against the supplied snapshot on 2026-10-11. The label is preserved for provenance and is not an independent verification of present open status or mathematical correctness.
\subsection{Short English statement}
Design reference material useful to intermediate users.
\subsection{Sources}
\begin{itemize}
\item[S1] Ulam AI and the UnsolvedMath Contributors, UnsolvedMath: A Curated Collection of Open Mathematics Problems, catalogue entry. Catalogue curation and metadata: CC BY 4.0. The source extract is retained; definitions, English summaries and qualifications are editorial additions.
\url{https://huggingface.co/datasets/ulamai/UnsolvedMath}.
\item[S2] AIM Workshop Problem Lists, source linked by the catalogue.
\url{https://github.com/PreTeXtBook/community/wiki/Moderated-Problem-Session}.
\end{itemize}
\paragraph{Preparation and review.}
The original source and UnsolvedMath contributors retain attribution for the source extract. Codex (GPT-6; exact serving revision not exposed) assisted with definitions, English summaries, scope notes and formatting. No new human mathematical authorship or independently verified solution is asserted. Catalogue identity checked against the complete supplied registry and current repository. Distinct upstream records, including repeated questions and broad or damaged extracts, are retained for complete coverage of the source CS category. A repeated record is not a new independent conjecture. Complete mathematical-equivalence and current-literature review has not been performed. Preparation and final review follow contribution prompts C and D, with source gaps disclosed. Structural and rendering checks do not establish a new theorem or open status.
\label{end:statement}

\section{Literature and context check}
The primary community-session entry was read. It is an implementation or documentation agenda, and its editorial catalogue title supplies no theorem specification. The following explicitly chosen model tests one route; no actual PreTeXt conversion, documentation release, or user study is claimed.
\section{Research attempt: A prerequisite reference that covers topics but not a usable route}
As a restricted documentation model, let a learner's state be a subset of finitely many established capabilities. An instruction has a prerequisite set $P$ and adds one capability $q$ when $P$ is contained in the current state. A reference may mention every capability yet provide no route from the learner's starting state. For instance, from the empty state, instructions $\{b\}\to a$ and $\{a\}\to b$ cover both topics but neither can start.

For a finite list of such monotone instructions and an initial capability set $S_0$, repeatedly add the conclusion of every instruction whose prerequisites are already present. This process terminates after at most the number of capabilities many strict additions. Every added capability is achievable: induction on insertion stages provides a valid instruction sequence. Conversely every capability achievable by a finite valid instruction sequence enters the closure: induction along that sequence shows each prerequisite was eventually present and its conclusion was therefore added. Thus the fixed point exactly characterizes reachability in this model.

The cycle example has closure empty. Adding one base instruction $\varnothing\to a$ gives a valid route to both $a$ and $b$. This is a concrete correction to the mistaken criterion that topic coverage alone suffices. But the monotonicity assumption matters: if an instruction changes software versions or invalidates earlier setup, capability sets may lose elements and this closure algorithm can claim a route that cannot be followed. An actual intermediate-user reference requires specified starting knowledge, versioned commands, and observed learning outcomes; the finite closure result does not supply them.
\section{Verification and adversarial audit}
The finite examples and implications are checked directly in the argument. No production PreTeXt build or conformance test was run. The proof applies only to the declared model, with its input and identity assumptions.
\paragraph{Audit corrections.} The declared instruction model is monotone; cyclic prerequisite lists and state-invalidating instructions are explicitly distinguished. Topic coverage is not called usability evidence.
\paragraph{External computation.} None; symbolic arguments and boundary checks only.
\section{Scope and remaining gaps}
The result is an auxiliary route analysis, not completion of the original software or documentation agenda. Real versions, supported syntax, transformations, and usability requirements need concrete implementation and testing. No mathematical novelty or current feature-completion claim is made.
\paragraph{Integration review.} Codex (GPT-6; exact serving revision not exposed) reviewed the written arguments and expanded layout and source-reference wrappers. Repository metadata and the production web renderer were checked using Python 3.13 and JavaScript with MathJax 3.2.2. This model-assisted review is not independent human proof verification; no PDF compilation or formal proof checking is asserted.

\section{Completion Estimate}
\noindent\textbf{COMPLETION: 0\%}
The calculations below the preserved source statement investigate a concrete route,
but establish no substantiated new progress toward the entire original target.
Known special cases and failed sufficient conditions are not credited as a solution.
\section{References}
The original source references and catalogue attribution are preserved in the
Sources section above. The following URLs were actually consulted for this attempt;
retention of an additional original reference does not claim fresh consultation of it.
\begin{itemize}
\item \url{https://github.com/PreTeXtBook/community/wiki/Moderated-Problem-Session} (consulted 2026-10-11). 
\end{itemize}
\section{Model and contributor disclosure}
OpenAI gpt-6-astra, requested reasoning configuration ultra, generated this research
attempt on 2026-10-11. The exact serving revision is not exposed. Codex GPT-6
performed orchestration and repository-format review; its exact serving revision
is also not exposed. The mathematical route, boundary checks and final adversarial
self-audit followed contribution prompts A then B in
\texttt{docs/contribution\_prompts.md}, after reading \texttt{CONTRIBUTING.md}.
Python 3.14.2 assembled the standalone documents; any mathematical computation is
specified in Verification. No theorem prover was used. No human mathematical
authorship, substantial human proof correction, or independent human proof checking
is asserted. Model-assisted review, source retrieval, and structural checks do not
establish correctness or novelty. Full private reasoning transcripts are unavailable
as reproducibility artifacts; the complete arguments, inputs and outcomes actually
used are stated here or in the companion batch audit.
\section{Final status}
Unresolved as a new research contribution. This is a concrete route investigation,
not a claim of new complete resolution or a current-status certification. A known
answer reproduced for context does not become a new contribution. The catalogue's
source-status label is preserved; this attempt status does not assert that a
source-labelled solved question is open.
\end{document}
